\documentclass[10pt,letterpaper]{article}
\usepackage[utf8]{inputenc}
\usepackage[intlimits]{amsmath}
\usepackage{amsfonts}
\usepackage{amssymb}
\usepackage{ragged2e}
\usepackage[letterpaper, margin=1in]{geometry}
\usepackage{graphicx}
\usepackage{cancel}
\usepackage{mathtools}
\usepackage{tabularx}
\usepackage{arydshln}
\usepackage{tensor}
\usepackage{array}
\usepackage{xcolor}
\usepackage[boxed]{algorithm}
\usepackage[noend]{algpseudocode}
\usepackage{listings}
\usepackage{textcomp}
% \usepackage[pdf,tmpdir,singlefile]{graphviz}
\usepackage{mathrsfs}
\usepackage{bbm}
\usepackage{tikz}
\usepackage{tikz-cd}
\usepackage{enumitem}
\usepackage{arydshln}
\usepackage{relsize}
\usepackage{multicol}
\usepackage{scalerel}

\usetikzlibrary{bayesnet}
\setlist{noitemsep}

%%%%%%%%%%%%%%%%%%%%%%%%%%%%%
% Formatting commands
%%%%%%%%%%%%%%%%%%%%%%%%%%%%%
\newcommand{\n}{\hfill\break}
\newcommand{\up}{\vspace{-\baselineskip}}
\newcommand{\hangblock}[2]{\par\noindent\settowidth{\hangindent}{\textbf{#1: }}\textbf{#1: }\!\!\!#2}
\newcommand{\lemma}[1]{\hangblock{Lemma}{#1}}
\newcommand{\defn}[1]{\hangblock{Defn}{#1}}
\newcommand{\thm}[1]{\hangblock{Thm}{#1}}
\newcommand{\cor}[1]{\hangblock{Cor}{#1}}
\newcommand{\prop}[1]{\hangblock{Prop}{#1}}
\newcommand{\ex}[1]{\hangblock{Ex}{#1}}
\newcommand{\exer}[1]{\hangblock{Exer}{#1}}
\newcommand{\fact}[1]{\hangblock{Fact}{#1}}
\newcommand{\remark}[1]{\hangblock{Remark}{#1}}
\newcommand{\proven}{\;$\square$\n}
\newcommand{\problem}[1]{\par\noindent{#1}\n}
\newcommand{\problempart}[2]{\par\noindent\indent{}\settowidth{\hangindent}{\textbf{(#1)} \indent{}}\textbf{(#1)} #2\n}
\newcommand{\ptxt}[1]{\textrm{\textnormal{#1}}}
\newcommand{\inlineeq}[1]{\centerline{$\displaystyle #1$}}
\newcommand{\pageline}{\noindent\rule{\textwidth}{0.1pt}}

%%%%%%%%%%%%%%%%%%%%%%%%%%%%%
% Math commands
%%%%%%%%%%%%%%%%%%%%%%%%%%%%%
% Set Theory
\newcommand{\card}[1]{\left|#1\right|}
\newcommand{\set}[1]{\left\{#1\right\}}
\newcommand{\setmid}{\;\middle|\;}
\newcommand{\ps}[1]{\mathcal{P}\left(#1\right)}
\newcommand{\pfinite}[1]{\mathcal{P}^{\ptxt{finite}}\left(#1\right)}
\newcommand{\naturals}{\mathbb{N}}
\newcommand{\N}{\naturals}
\newcommand{\integers}{\mathbb{Z}}
\newcommand{\Z}{\integers}
\newcommand{\rationals}{\mathbb{Q}}
\newcommand{\Q}{\rationals}
\newcommand{\reals}{\mathbb{R}}
\newcommand{\R}{\reals}
\newcommand{\complex}{\mathbb{C}}
\newcommand{\C}{\complex}
\newcommand{\halfPlane}{\mathbb{H}}
\let\H\relax
\newcommand{\H}{\halfPlane}
\newcommand{\comp}{^{\complement}}
\DeclareMathOperator{\Hom}{Hom}
\newcommand{\Ind}{\mathbbm{1}}
\newcommand{\cut}{\setminus}
\DeclareMathOperator{\elem}{elem}

% Graph Theory
\let\deg\relax
\DeclareMathOperator{\deg}{deg}
\newcommand{\degp}{\ptxt{deg}^{+}}
\newcommand{\degn}{\ptxt{deg}^{-}}
\newcommand{\precdot}{\mathrel{\ooalign{$\prec$\cr\hidewidth\hbox{$\cdot\mkern0.5mu$}\cr}}}
\newcommand{\succdot}{\mathrel{\ooalign{$\cdot\mkern0.5mu$\cr\hidewidth\hbox{$\succ$}\cr\phantom{$\succ$}}}}
\DeclareMathOperator{\cl}{cl}
\DeclareMathOperator{\affdim}{affdim}

% Probability
\newcommand{\parSymbol}{\P}
\newcommand{\Prob}{\mathbb{P}}
\renewcommand{\P}{\Prob}
\newcommand{\Avg}{\mathbb{E}}
\newcommand{\E}{\Avg}
\DeclareMathOperator{\Var}{Var}
\DeclareMathOperator{\cov}{cov}
\DeclareMathOperator{\Unif}{Unif}
\DeclareMathOperator{\Binom}{Binom}
\newcommand{\CI}{\mathrel{\text{\scalebox{1.07}{$\perp\mkern-10mu\perp$}}}}

% Standard Math
\newcommand{\inv}{^{-1}}
\newcommand{\abs}[1]{\left|#1\right|}
\newcommand{\ceil}[1]{\left\lceil{}#1\right\rceil{}}
\newcommand{\floor}[1]{\left\lfloor{}#1\right\rfloor{}}
\newcommand{\conj}[1]{\overline{#1}}
\newcommand{\of}{\circ}
\newcommand{\tri}{\triangle}
\newcommand{\inj}{\hookrightarrow}
\newcommand{\surj}{\twoheadrightarrow}
\newcommand{\ndiv}{\nmid}
\renewcommand{\epsilon}{\varepsilon}
\newcommand{\divides}{\mid}
\newcommand{\ndivides}{\nmid}
\DeclareMathOperator{\lcm}{lcm}
\DeclareMathOperator{\sgn}{sgn}
\newcommand{\map}[4]{\!\!\!\begin{array}[t]{rcl}#1 & \!\!\!\!\to & \!\!\!\!#2\\ #3 & \!\!\!\!\mapsto & \!\!\!\!#4\end{array}}
\newcommand{\bigsum}[2]{\smashoperator[lr]{\sum_{\scalebox{#1}{$#2$}}}}

% Linear Algebra
\newcommand{\Id}{\textrm{\textnormal{Id}}}
\newcommand{\im}{\textrm{\textnormal{im}}}
\newcommand{\norm}[1]{\abs{\abs{#1}}}
\newcommand{\tpose}{^{T}\!}
\newcommand{\iprod}[1]{\left<#1\right>}
\DeclareMathOperator{\tr}{tr}
\DeclareMathOperator{\trace}{tr}
\newcommand{\chgBasMat}[3]{\!\!\tensor*[_{#1}]{\left[#2\right]}{_{#3}}}
\newcommand{\vecBas}[2]{\tensor*[]{\left[#1\right]}{_{#2}}}
\DeclareMathOperator{\GL}{GL}
\DeclareMathOperator{\Mat}{Mat}
\DeclareMathOperator{\vspan}{span}
\DeclareMathOperator{\rank}{rank}
\newcommand{\V}[1]{\vec{#1}}
\DeclareMathOperator{\proj}{proj}
\DeclareMathOperator{\compProj}{comp}
\DeclareMathOperator{\row}{row}
\newcommand{\smallPMatrix}[1]{\paren{\begin{smallmatrix}#1\end{smallmatrix}}}
\newcommand{\smallBMatrix}[1]{\brack{\begin{smallmatrix}#1\end{smallmatrix}}}
\newcommand{\pmat}[1]{\begin{pmatrix}#1\end{pmatrix}}
\newcommand{\bmat}[1]{\begin{bmatrix}#1\end{bmatrix}}

% Multilinear Algebra
\newcommand{\Lsym}{\L}
\let\L\relax
\DeclareMathOperator{\L}{\mathscr{L}}
\DeclareMathOperator{\A}{\mathcal{A}}
\DeclareMathOperator{\Alt}{Alt}
\DeclareMathOperator{\Sym}{Sym}
\newcommand{\ot}{\otimes}
\newcommand{\ox}{\otimes}
\DeclareMathOperator{\asc}{asc}
\DeclareMathOperator{\asSet}{set}
\DeclareMathOperator{\sort}{sort}
\DeclareMathOperator{\ringA}{\mathring{A}}

% Topology
\newcommand{\closure}[1]{\overline{#1}}
\newcommand{\uball}{\mathcal{U}}
\DeclareMathOperator{\Int}{Int}
\DeclareMathOperator{\Ext}{Ext}
\DeclareMathOperator{\Bd}{Bd}
\DeclareMathOperator{\rInt}{rInt}
\DeclareMathOperator{\ch}{ch}
\DeclareMathOperator{\ah}{ah}
\newcommand{\LargerTau}{\mathlarger{\mathlarger{\mathlarger{\mathlarger{\tau}}}}}
\newcommand{\Tau}{\mathcal{T}}

% Analysis
\DeclareMathOperator{\Graph}{Graph}
\DeclareMathOperator{\epi}{epi}
\DeclareMathOperator{\hypo}{hypo}
\DeclareMathOperator{\supp}{supp}
\newcommand{\lint}[2]{\underset{#1}{\overset{#2}{{\color{black}\underline{{\color{white}\overline{{\color{black}\int}}\color{black}}}}}}}
\newcommand{\uint}[2]{\underset{#1}{\overset{#2}{{\color{white}\underline{{\color{black}\overline{{\color{black}\int}}\color{black}}}}}}}
\newcommand{\alignint}[2]{\underset{#1}{\overset{#2}{{\color{white}\underline{{\color{white}\overline{{\color{black}\int}}\color{black}}}}}}}
\newcommand{\extint}{\ptxt{ext}\int}
\newcommand{\extalignint}[2]{\ptxt{ext}\alignint{#1}{#2}}
\newcommand{\conv}{\ast}
\newcommand{\pd}[2]{\frac{\partial{}#1}{\partial{}#2}}
\newcommand{\del}{\nabla}
\DeclareMathOperator{\grad}{grad}
\DeclareMathOperator{\curl}{curl}
\let\div\relax
\DeclareMathOperator{\div}{div}
\DeclareMathOperator{\vol}{vol}

% Complex Analysis
\let\Re\relax
\DeclareMathOperator{\Re}{Re}
\let\Im\relax
\DeclareMathOperator{\Im}{Im}
\DeclareMathOperator{\Res}{Res}

% Abstract Algebra
\DeclareMathOperator{\ord}{ord}
\newcommand{\generated}[1]{\left<#1\right>}
\newcommand{\cycle}[1]{\smallPMatrix{#1}}
\newcommand{\id}{\ptxt{id}}
\newcommand{\iso}{\cong}
\DeclareMathOperator{\Aut}{Aut}
\DeclareMathOperator{\SL}{SL}
\DeclareMathOperator{\op}{op}
\newcommand{\isom}[4]{\!\!\!\begin{array}[t]{rcl}#1 & \!\!\!\!\overset{\sim}{\to} & \!\!\!\!#2\\ #3 & \!\!\!\!\mapsto & \!\!\!\!#4\end{array}}
\newcommand{\F}{\mathbb{F}}
\newcommand{\acts}{\;\rotatebox[origin=c]{-90}{$\circlearrowright$}\;}
\newcommand{\disjunion}{\mathrel{\text{\scalebox{1.07}{$\perp\mkern-10mu\perp$}}}}
\DeclareMathOperator{\SO}{SO}
\DeclareMathOperator{\stab}{stab}
\DeclareMathOperator{\nullity}{nullity}
\DeclareMathOperator{\Perm}{Perm}
\DeclareMathOperator{\nsubgp}{\vartriangleleft}
\DeclareMathOperator{\notnsubgp}{\ntriangleleft}

% Convex Optimization
\newcommand{\sectionSymbol}{\S}
\let\S\relax
\newcommand{\S}{\mathbb{S}}
\DeclareMathOperator{\dist}{dist}
\DeclareMathOperator{\dom}{dom}
\DeclareMathOperator{\diag}{diag}
\DeclareMathOperator{\ones}{\mathbbm{1}}
\newcommand{\minimizeOver}[3]{\begin{array}{rl}\underset{#1}{\ptxt{minimize}} & #2\\ \ptxt{subject to} & #3\end{array}}
\newcommand{\maximizeOver}[3]{\begin{array}{rl}\underset{#1}{\ptxt{maximize}} & #2\\ \ptxt{subject to} & #3\end{array}}
\newcommand{\minimizationProblem}[2]{\minimizeOver{}{#1}{#2}}
\newcommand{\maximizationProblem}[2]{\maximizeOver{}{#1}{#2}}

% Proofs
\newcommand{\st}{s.t.}
\newcommand{\unique}{!}
\newcommand{\iffdef}{\overset{\ptxt{def}}{\Leftrightarrow}}
\newcommand{\eqdef}{\overset{\ptxt{def}}{=}}
\newcommand{\eqVertical}{\rotatebox[origin=c]{90}{=}}
\newcommand{\mapsfrom}{\mathrel{\reflectbox{\ensuremath{\mapsto}}}}
\newcommand{\mapsdown}{\rotatebox[origin=c]{-90}{$\mapsto$}\mkern2mu}
\newcommand{\mapsup}{\rotatebox[origin=c]{90}{$\mapsto$}\mkern2mu}
\newcommand{\from}{\!\mathrel{\reflectbox{\ensuremath{\to}}}}

% Brackets
\newcommand{\paren}[1]{\left(#1\right)}
\renewcommand{\brack}[1]{\left[#1\right]}
\renewcommand{\brace}[1]{\left\{#1\right\}}
\newcommand{\ang}[1]{\left<#1\right>}

% Algorithms
\algrenewcommand{\algorithmiccomment}[1]{\hskip 1em \texttt{// #1}}
\algrenewcommand\algorithmicrequire{\textbf{Input:}}
\algrenewcommand\algorithmicensure{\textbf{Output:}}
\newcommand{\algP}{\ptxt{\textbf{P}}}
\newcommand{\algNP}{\ptxt{\textbf{NP}}}
\newcommand{\algNPC}{\ptxt{\textbf{NP-Complete}}}
\newcommand{\algNPH}{\ptxt{\textbf{NP-Hard}}}
\newcommand{\algEXP}{\ptxt{\textbf{EXP}}}

%%%%%%%%%%%%%%%%%%%%%%%%%%%%%
% Other commands
%%%%%%%%%%%%%%%%%%%%%%%%%%%%%
\newcommand{\flag}[1]{\textbf{\textcolor{red}{#1}}}
\newcommand{\uSym}{\u}
\let\u\relax
\newcommand{\u}[1]{\underline{#1}}
\newcommand{\bSym}{\b}
\let\b\relax
\newcommand{\b}[1]{\textbf{#1}}
\newcommand{\iSym}{\i}
\let\i\relax
\newcommand{\i}[1]{\textit{#1}}

%%%%%%%%%%%%%%%%%%%%%%%%%%%%%%%%%%%%%%%
% Make l's curvy in math environments %
%%%%%%%%%%%%%%%%%%%%%%%%%%%%%%%%%%%%%%%
\mathcode`l="8000
\begingroup
\makeatletter
\lccode`\~=`\l
\DeclareMathSymbol{\lsb@l}{\mathalpha}{letters}{`l}
\lowercase{\gdef~{\ifnum\the\mathgroup=\m@ne \ell \else \lsb@l \fi}}%
\endgroup

%%%%%%%%%%%%%%%%%%%%%%%%%
% Fix \vdots and \ddots %
%%%%%%%%%%%%%%%%%%%%%%%%%
\usepackage{letltxmacro}
\LetLtxMacro\orgvdots\vdots
\LetLtxMacro\orgddots\ddots

\makeatletter
\DeclareRobustCommand\vdots{%
	\mathpalette\@vdots{}%
}
\newcommand*{\@vdots}[2]{%
	% #1: math style
	% #2: unused
	\sbox0{$#1\cdotp\cdotp\cdotp\m@th$}%
	\sbox2{$#1.\m@th$}%
	\vbox{%
		\dimen@=\wd0 %
		\advance\dimen@ -3\ht2 %
		\kern.5\dimen@
		% remove side bearings
		\dimen@=\wd2 %
		\advance\dimen@ -\ht2 %
		\dimen2=\wd0 %
		\advance\dimen2 -\dimen@
		\vbox to \dimen2{%
			\offinterlineskip
			\copy2 \vfill\copy2 \vfill\copy2 %
		}%
	}%
}
\DeclareRobustCommand\ddots{%
	\mathinner{%
		\mathpalette\@ddots{}%
		\mkern\thinmuskip
	}%
}
\newcommand*{\@ddots}[2]{%
	% #1: math style
	% #2: unused
	\sbox0{$#1\cdotp\cdotp\cdotp\m@th$}%
	\sbox2{$#1.\m@th$}%
	\vbox{%
		\dimen@=\wd0 %
		\advance\dimen@ -3\ht2 %
		\kern.5\dimen@
		% remove side bearings
		\dimen@=\wd2 %
		\advance\dimen@ -\ht2 %
		\dimen2=\wd0 %
		\advance\dimen2 -\dimen@
		\vbox to \dimen2{%
			\offinterlineskip
			\hbox{$#1\mathpunct{.}\m@th$}%
			\vfill
			\hbox{$#1\mathpunct{\kern\wd2}\mathpunct{.}\m@th$}%
			\vfill
			\hbox{$#1\mathpunct{\kern\wd2}\mathpunct{\kern\wd2}\mathpunct{.}\m@th$}%
		}%
	}%
}
\makeatother

\newcommand{\B}{
	\begin{tikzpicture}
	\filldraw [fill=red, draw=black] (0, 0) rectangle (0.37, 0.45);
	\draw [line width=0.5mm, white ] (0.1,0.08) -- (0.1,0.38);
	\draw[line width=0.5mm, white ] (0.1, 0.35) .. controls (0.2, 0.35) and (0.4, 0.2625) .. (0.1, 0.225);
	\draw[line width=0.5mm, white ] (0.1, 0.225) .. controls (0.2, 0.225) and (0.4, 0.1625) .. (0.1, 0.1);
	\end{tikzpicture}
}

\author{Thomas Cohn}
\title{Math 493 Lecture 15}
\date{10/28/19} % Can also use \today

\begin{document}
\maketitle
\setlength\RaggedRightParindent{\parindent}
\RaggedRight

\par\noindent
Let $G$ be a group.\n

\par\noindent
Recall: for $A,B\subset{}G$ subsets, $AB=\set{ab\mid{}a\in{}A,b\in{}B}$. If $A$ and $B$ are subgroups, then $AB$ is not always a subgroup. However, if $A$ is normal, then $AB$ is a subgroup.\n

\par\noindent
Recall:
\defn{
	$G$ is the \textbf{(internal) direct product} of subgroups $A$ and $B$ if
	\begin{enumerate}[leftmargin=4\parindent]
		\item $A$ and $B$ are normal
		\item $AB=G$
		\item $A\cap{}B=\set{1}$
	\end{enumerate}\up\n
	Notation: $G=A\times{}B$\n
}

\defn{
	$G$ is the \textbf{(internal) semi-direct product} of subgroups $A$ and $B$ if
	\begin{enumerate}[leftmargin=4\parindent]
		\item $A$ is normal
		\item $AB=G$
		\item $A\cap{}B=\set{1}$
	\end{enumerate}\up\n
	Notation: $G=A\rtimes{}B$\n
}

\ex{
	\begin{enumerate}[leftmargin=4\parindent]
		\item $G=M$, the group of rigid motions of the plane $P=\R^{2}$.\n
		$A$ is the group of translations.\n
		$B$ is $\operatorname{O}(2)$, the group of origin-preserving transformations.\n
		$G=A\rtimes{}B$.\n
		\item $G=D_{n}$, the dihedral group of order $2n$.\n
		$A$ is the group of rotations ($\ord(A)=n$).\n
		$B$ is the subgroup generated by any reflection.\n
		$G=A\rtimes{}B$.\n
		\item $G=\operatorname{O}(n)$.\n
		$A=\SO(n)$.\n
		$B$ is the subgroup generated by any reflection.\n
		$G=A\rtimes{}B$.\n
		\item $G=A\times{}B$.\n
		$A=A$.\n
		$B=B$.\n
		$G=A\rtimes{}B$. (I.e. every direct product is a semi-direct product.)
	\end{enumerate}
}

\par\noindent
So, among other things, we have $\Z/n\Z\times\Z/2\Z=\Z/n\Z\rtimes\Z/2\Z=D_{n}$. Thus, we cannot recover $G$ from $A$ and $B$ as an abstract group from a semi-direct product.\n

\par\noindent
Let $X$ be a set. $\Perm(X)$ is the group of permutations of $X$, i.e., the set of all bijections $X\to{}X$, under composition.\n
$S_{n}=\Perm(\set{1,\ldots,n})$.\n

\par\noindent
If $X$ is a $G$-set, then for $g\in{}G$, we get a permutation of $X$ by $\map{X}{X}{x}{gx}$.\n
This defines a function $G\to\Perm(X)$, which is a group homomorphism. In fact, giving an action of $G$ on $X$ is the same as giving a homomorphism $G\to\Perm(X)$.\n

\defn{
	Now, say $X$ is a group. $\Aut(X)\subseteq\Perm(X)$. If $G$ acts on $X$, we say $G$ \textbf{acts by group homomorphisms} if the map $G\to\Perm(X)$ lands in $\Aut(X)$.\n
}

\par\noindent
Explicitly, this means $g\in{}G,x,y\in{}X\Rightarrow{}g\cdot(xy)=(g\cdot{}x)(g\cdot{}y)$.\n

\ex{
	\begin{enumerate}[leftmargin=4\parindent]
		\item $G$ acts on itself by left-multiplication. This action \textit{is not} by group homomorphisms.
		\item $G$ acts on itself by conjugation. This action \textit{is} by group homomorphisms.
	\end{enumerate}
}

\par\noindent
Say $G=A\rtimes{}B$. Given $b\in{}B$, we know that conjugation by $b$ is a group homomorphism $G\to{}G$, and maps $A$ to itself because $A$ is normal. Let
\[
	\map{\varphi_{b}:A}{A}{a}{bab\inv}
\]
Then $\varphi_{b}\in\Aut(A)$. The function
\[
	\map{\varphi:B}{\Aut(A)}{b}{\varphi_{b}}
\]
is a group homomorphism. Reason: $\varphi_{bb'}(a)=bb'a(bb')\inv=b(b'a(b')\inv)b\inv=\varphi_{b}(\varphi_{b'}(a))$.\n

\par\noindent
Note: Every element of $G$ can be written \textit{uniquely} as $ab$ with $a\in{}A$ and $b\in{}B$.\n
Existence is clear -- we just need to check uniqueness.\n
Say $ab=a'b'$. Then $1=(a')\inv{}a=b'b\inv\in{}A\cap{}B=\set{1}$. So $a=a'$ and $b=b'$, because inverses are unique.\n

\par\noindent
To rephrase, the function $\map{A\times{}B}{G}{(a,b)}{ab}$ is a bijection of sets. But in general, it \textit{is not} a group homomorphism.\n

\par\noindent
How do we multiply elements of $G$ under this description?
\[
	(ab)(a'b')=(aba'b\inv)(bb')=(a\varphi_{b}(a'))(bb')
\]

\defn{
	$A$ and $B$ are two groups. $\varphi:B\to\Aut(A)$ is a group homomorphism. The \textbf{(external) semi-direct product} of $A$ and $B$ is the set of elements in $A\times{}B$, with multiplication\n
	$(a,b)(a',b')=(a\varphi_{b}(a'),bb')$. Notation: $G=A\rtimes_{\varphi}B$ (the subscript is optional).\n
}

\par\noindent
Suppose $G=A\rtimes{}B$. We've constructed $\varphi:B\to\Aut(A)$.\n

\par\noindent
From the above discussion, we know $G\cong{}A\rtimes_{\varphi}B$ (external). So internal semi-direct products are external semi-direct products.\n

\par\noindent
Now, say $G=A\rtimes_{\varphi}B$ (external). Then $\bar{A}=\set{(a,1)\mid{}a\in{}A}$, $\bar{B}=\set{(1,b)\mid{}b\in{}B}$. $\bar{A},\bar{B}\subseteq{}G$. We claim $\bar{A}$ and $\bar{B}$ are subgroups of $G$, and $G$ is the internal semi-direct product.\n
Proof:\n
$\bar{A}$ is closed under multiplication: $(a,1)(a',1)=(a\varphi_{1}(a'),1\cdot{}1)=(aa',1)$, because $\varphi_{1}=\Id_{A}$, so $\bar{A}\cong{}A$.\n
$\bar{B}$ is closed under multiplication: $(1,b)(1,b')=(1\varphi_{b}(1),bb')=(1,bb')$, so $\bar{B}\cong{}B$.\n
$G=\bar{A}\bar{B}$: $(a,1)(1,b)=(a\varphi_{1}(1),1\cdot{}b)=(a,b)$.\n
$\bar{A}$ normal: $(1,b)(a,1)(1,b)\inv=(1,b)(a,1)(1,b\inv)=(1,b)(a,b\inv)=(\varphi_{b}(a),bb\inv)=(\varphi_{b}(a),1)\in\bar{A}$.\proven

\ex{
	Some external semi-direct products:
	\begin{enumerate}[leftmargin=4\parindent]
		\item $\varphi:B\to\Aut(A)$ is the trivial homomorphism. Then $\varphi_{b}=\Id_{A}$, $\forall{}b$, so $\varphi_{b}(a)=a$, $\forall{}a,b$.\n
		Then $(a,b)(a',b')=(a\varphi_{b}(a'),bb')=(aa',bb')$. So we get the direct product $A\times{}B$.\n
		\item $G=D_{n}$, $A$ is the group of rotations $\generated{\rho_{\frac{2\pi}{n}}}\cong\Z/n\Z$. $B$ is the group generated by some reflection, which is equal to $\generated{r}\cong\Z/2\Z$. We want to understand
		\[
			\map{\varphi:B}{\Aut(A)}{1}{\Id_{A}\\ r & \!\!\!\!\mapsto & \!\!\!\!(\rho^{k}\mapsto{}\rho^{-k})}
		\]
		Recall, $r\rho^{k}r\inv=\rho^{-k}$, so we have
		\begin{center}
			\begin{tikzcd}[ampersand replacement=\&]
				\&[-25pt] \rho^{k} \& k \arrow[l,mapsto]{}\\[-15pt]
				\rho^{k} \arrow[d,mapsto]{} \& \generated{\rho} \arrow{d}{\varphi_{r}} \& \Z/n\Z \arrow{l}[swap]{\sim} \arrow{d} \&[-25pt] k \arrow[d,mapsto]{}\\
				\rho^{-k} \& \generated{\rho} \& \Z/n\Z \arrow{l}[swap]{\sim} \& -k\\
			\end{tikzcd}
		\end{center}
		\item $F$ a field, $A=F^{n}$ (the group of column vectors, under addition). $B=\GL_{n}(F)\subseteq\Aut(A)$. ($B$ consists of $F$-linear automorphisms.)Take $\varphi$ to be inclusion. Explicitly, $\varphi_{b}(a)=ba$ (matrix multiplication).\n
		$G=A\rtimes_{\varphi}B$: $(a,b)(a',b')=(a+\varphi_{b}(a'),bb')=(a+ba',bb')$ (with vector addition, and matrix multiplication).
	\end{enumerate}
}

\par\noindent
Let $A\nsubgp{}G$ a normal subgroup. Let $A'$ be some nonzero proper $F$-subspace of $A$. $A'\nsubgp{}A$, but $A'\notnsubgp{}G$.\n

\par\noindent
Question: What is $\Aut(\Z/n\Z)$?\n
Suppose $\varphi:\Z/n\Z\to\Z/n\Z$ is an automorphism. $\varphi$ is determined by $\varphi(1)$, call this $a\in\Z/n\Z$.\n
$\varphi(2)=\varphi(1+1)=\varphi(1)+\varphi(1)=a+a=2a$.\n
$\varphi(3)=\varphi(2+1)=\varphi(2)+\varphi(1)=2a+a=3a$.\n
In general, $\varphi(k)=ka$ (multiplication modulo $n$). Since $\varphi$ is an automorphism, $\exists{}k$ \st{} $\varphi(k)=1$, so $ka=1$. So $a$ is invertible under multiplication.\n

\defn{
	We say $a$ is a \textbf{unit} of $\Z/n\Z$.\n
}

\par\noindent
$(\Z/n\Z)^{\times}$ is the set of units of $\Z/n\Z$, and it is a group under multiplication.\n

\par\noindent
Summary: The map $\map{(\Z/n\Z)^{\times}}{\Aut(\Z/n\Z)}{a}{m_{a}=\ptxt{multiplication by }a}$ is an isomorphism of groups.\n
Note: $m_{a}(m_{b}(x))=a\cdot{}b\cdot{}x=m_{ab}(x)$, so $m_{a}\of{}m_{b}=m_{ab}$, so this is a group homomorphism.)\n

\lemma{
	If $a\in\Z/n\Z$, then $a$ is a unit iff $\gcd(a,n)=1$.\n
}

\newpage
\ex{
	\begin{itemize}[leftmargin=4\parindent]
		\item If $n=p$ is prime, $(\Z/n\Z)^{\times}=\F_{p}^{\times}=\F\cut\set{0}$. We've shown this is cyclic, and of order $p-1$.
		\item $(\Z/4\Z)^{\times}=\set{1,3}\cong\Z/2\Z$.
		\item $(\Z/6\Z)^{\times}=\set{1,5}\cong\Z/2\Z$.
		\item $(\Z/8\Z)^{\times}=\set{1,3,5,7}\cong\Z/2\Z\times\Z/2\Z$.
		\item $(\Z/16\Z)^{\times}\cong\Z/2\Z\times\Z/4\Z$.
	\end{itemize}
}

\ex{
	The last computation implies we have an injective group homomorphism $\varphi:\Z/4\Z\to\Aut(\Z/16\Z)$. We get a semi-direct product $\Z/16\Z\rtimes_{\varphi}\Z/4\Z$.\n
}

\end{document}